% 来源B：提取自教材帮(1).pdf，PDF第1页，书页100。
% 本分块保留知识点3、知识点4及例5、例6。
% 目标：运算和复合的单调性及适用前提；方法：比较函数值差的符号；
% 易错点：跨零点取倒数、内层越出外层定义域；难度：基础到中档；评价：口述条件并解释例6。

\begin{frame}[label=Doperations]
  \PageHeading{补充知识\quad 单调函数的运算性质}
  设 $f(x),g(x)$ 在同一区间 $I$ 上单调，“增、减”指严格单调。
  \par\vspace{0.2cm}
  \begin{enumerate}
    \item $f(x)$ 与 $f(x)+C$ 的单调性相同，$C$ 为常数。
    \item 当 $a>0$ 时，$af(x)$ 与 $f(x)$ 的单调性相同；\par
          当 $a<0$ 时，两者的单调性相反。
    \item 当 $f(x)\ge0$ 时，$\sqrt{f(x)}$ 与 $f(x)$ 的单调性相同。
  \end{enumerate}
  \pause
  \Analysis{加常数不改变函数值的差；数乘取决于系数符号；\par
  平方根函数在 $[0,+\infty)$ 上严格递增。$a=0$ 时得到常函数，另行处理。}
\end{frame}

\begin{frame}[label=Dsum]
  \PageHeading{补充知识\quad 和与差的单调性}
  在 $f(x),g(x)$ 的共同单调区间上：
  \begin{center}
  \renewcommand{\arraystretch}{1.5}
  \begin{tabular}{c|c|c|c}
    $f(x)$ & $g(x)$ & $f(x)+g(x)$ & $f(x)-g(x)$\\\hline
    增 & 增 & 增 & 不能确定\\
    增 & 减 & 不能确定 & 增\\
    减 & 减 & 减 & 不能确定\\
    减 & 增 & 不能确定 & 减
  \end{tabular}
  \end{center}
  \pause
  \Analysis{将减法看作加上 $-g(x)$；“不能确定”表示仅凭两者的单调性不足以判断。
  例如 $2x-x=x$ 递增，而 $x-2x=-x$ 递减。}
\end{frame}

\begin{frame}[label=D05q]
  \PageHeading{补充例5\quad 倒数与单调性}
  证明：若 $f(x)$ 在区间 $I$ 上单调，且恒为正值或恒为负值，\par
  $f(x)$ 与 $\dfrac1{f(x)}$ 具有相反的单调性。
  \pause
  \Analysis{设 $x_1,x_2\in I$，且 $x_1<x_2$，则
  \[
  \frac1{f(x_2)}-\frac1{f(x_1)}
   =-\frac{f(x_2)-f(x_1)}{f(x_1)f(x_2)}.
  \]
  因为 $f(x_1)f(x_2)>0$，倒数之差与原函数值之差异号。\par
  因而原函数递增时，倒数函数递减；原函数递减时，倒数函数递增。}
\end{frame}

\begin{frame}[label=Dreciprocal]
  \PageHeading{补充知识\quad 倒数与乘积的条件}
  \textbf{倒数：}当 $f(x)$ 单调且恒正或恒负，$a$ 为非零常数时，
  \[
  a>0:\quad \frac a{f(x)}\text{ 与 }f(x)\text{ 单调性相反};\qquad
  a<0:\quad \text{单调性相同}.
  \]
  \pause
  \textbf{乘积：}若 $f(x),g(x)$ 的单调性相同，则
  \begin{center}
  \renewcommand{\arraystretch}{1.3}
  \begin{tabular}{c|c|c}
    $f,g$ 的符号 & $f,g$ 都递增 & $f,g$ 都递减\\\hline
    都不小于零 & $fg$ 递增 & $fg$ 递减\\
    都不大于零 & $fg$ 递减 & $fg$ 递增
  \end{tabular}
  \end{center}
  \pause
  \par\vspace{0.1cm}
  \textbf{注意：}相乘的规则必须结合函数值的符号，不能只看增减。
\end{frame}

\begin{frame}[label=Dcomposition]
  \PageHeading{补充知识\quad 复合函数的单调性}
  设 $g(x)$ 在区间 $I$ 上单调，且外层函数 $f(t)$ 在内层值域 $g(I)$ 上单调。
  \begin{center}
  \renewcommand{\arraystretch}{1.4}
  \begin{tabular}{c|c|c}
    内层 $t=g(x)$ & 外层 $y=f(t)$ & $y=f(g(x))$\\\hline
    增 & 增 & 增\\
    减 & 减 & 增\\
    减 & 增 & 减\\
    增 & 减 & 减
  \end{tabular}
  \end{center}
  \pause
  \textbf{简记：}同增异减。\quad 先确认复合有意义，再使用增减规则。
\end{frame}

\begin{frame}[label=Dsteps]
  \PageHeading{复合函数的判断步骤}
  \begin{enumerate}
    \item 将复合函数分解为内层函数和外层函数。
    \item 确定复合函数的定义域。
    \item 分区间判断内层单调性，并求该区间上的内层取值范围。
    \item 在对应取值范围内判断外层的单调性，再合成结论。
  \end{enumerate}
  \pause
  \par\vspace{0.2cm}
  \textbf{关键：}外层函数所用区间，要与内层在讨论区间上的值域对应。
\end{frame}

\begin{frame}[label=D06q]
  \PageHeading{补充例6\quad 复合函数的单调区间}
  已知函数 $f(x)$ 在定义域 $[0,+\infty)$ 上单调递减，\par
  则 $f(1-x^2)$ 的单调递减区间是\underline{\hspace{2cm}}。
  \pause
  \Answer{$[-1,0]$。}
  \Analysis{由 $1-x^2\ge0$ 得 $-1\le x\le1$。令 $u=1-x^2$，则
  \begin{center}
  \renewcommand{\arraystretch}{1.3}
  \begin{tabular}{c|c|c|c}
    $x$ 的区间 & 内层 $u$ & 外层 $f(u)$ & 复合函数\\\hline
    $[-1,0]$ & 增 & 减 & 减\\
    $[0,1]$ & 减 & 减 & 增
  \end{tabular}
  \end{center}}
\end{frame}
